\SourceLeaf{166}
\begin{center}
{\large POSTHUMOUS PAPERS AND CORRESPONDENCE.\par}
\medskip
[ON THE FIRST PRINCIPLES OF GEOMETRY.]
\end{center}

\bigskip
\begin{center}\textsc{Gauss} to \textsc{Bolyai}. Helmstedt, 16 December 1799.\end{center}

\medskip
.\, .\, .\, .\, .\, .\, . I very much regret that I did not take advantage of our former proximity to learn more about your work on the first principles of geometry. It would certainly have spared me many fruitless efforts and made me more at ease than someone like me can be while there is still so much to be desired in such a subject. I have made considerable progress in my own work on it (although my other, entirely unrelated duties leave me little time for it). But the path I have taken leads less to the goal one desires, and which you assure me you have attained, than to calling the truth of geometry into question. I have indeed come across some things that most people would already count as a proof, but that in my eyes prove practically NOTHING. For example, if it could be proved that a rectilinear triangle with an area greater than any given area is possible, I could prove the whole of geometry with complete rigor. Most people would probably accept that as an axiom. I would not. It is quite possible that no matter how far apart the three vertices of the triangle are taken in space, the area would still always lie below (infra) a given bound. I have several propositions of this kind, but I find none of them satisfactory. Please do make
